\section{Процедура HandleCmdImportShow}
\footnotesize
\begin{listing}{1}
/*************************************************************************
* Процедура     : HandleCmdImportShow                                    *
*------------------------------------------------------------------------*
* Назначение    : Обработка команд импорта и отображения данных          *
* Входные       : hWnd - дескриптор окна, wmId - идентификатор команды   *
* Выходные      : Нет                                                    *
*------------------------------------------------------------------------*
* ОПИСАНИЕ ФУНКЦИОНИРОВАНИЯ:                                             *
* Обрабатывает нажатия кнопок импорта и показа данных, управляет         *
* состоянием интерфейса и сбросом предыдущих результатов.                *
*************************************************************************/
void HandleCmdImportShow(HWND hWnd, int wmId) {
    if (wmId == ID_BTN_IMPORT) { wstring f;
        if (GetLoadFileNameExplorer(hWnd, f)) {
            if (ImportDataFromFile(f)) {
                g_StatusMessage = L"✔ Успех! Вариант №" + to_wstring(g_VariantNum);
                g_ActiveView = 1; g_IsCalculated = false;
                if (g_hPlotBitmapA) { DeleteObject(g_hPlotBitmapA); g_hPlotBitmapA = NULL;}
                if (g_hPlotBitmapB) { DeleteObject(g_hPlotBitmapB); 
		g_hPlotBitmapB = NULL; }
	} else ShowError(hWnd, ErrParse);}
}
    else if (wmId == ID_BTN_SHOW) { if (g_VariantNum == 0) ShowError(hWnd, WarnNoData);
        else { g_StatusMessage = L"► Исходные данные ВАР №" + to_wstring(g_VariantNum); 
            g_ActiveView = 1; }
    }
}
\end{listing}
\normalsize

\pagebreak

\section{Процедура HandleCmdCalcExport}
\footnotesize
\begin{listing}{1}
/*************************************************************************
* Процедура     : HandleCmdCalcExport                                    *
*------------------------------------------------------------------------*
* Назначение    : Обработка команд расчета и экспорта                    *
* Входные       : hWnd - дескриптор окна, wmId - идентификатор команды   *
* Выходные      : Нет                                                    *
*------------------------------------------------------------------------*
* ОПИСАНИЕ ФУНКЦИОНИРОВАНИЯ:                                             *
* Выполняет расчет параметров поля и экспорт отчета в файл.              *
*************************************************************************/
void HandleCmdCalcExport(HWND hWnd, int wmId) {
    if (wmId == ID_BTN_CALC) {
        if (g_VariantNum == 0) ShowError(hWnd, WarnNoData);
        else { g_StatusMessage = L"✔ Расчёт выполнен!"; g_ActiveView = 2; 
            g_IsCalculated = true;}
    }
    else if (wmId == ID_BTN_EXPORT) {
        if (!g_IsCalculated) { ShowError(hWnd, WarnNoCalc); 
            return;}
        wstring s;
        if (GetSaveFileNameExplorer(hWnd, s)) {
            if (ExportDataToFile(s)) 
                g_StatusMessage = L"✔ Отчёт сохранён!";
            else 
                ShowError(hWnd, ErrWrite);}
    }
}
\end{listing}
\normalsize

\pagebreak

\section{Процедура HandleCmdVis}
\footnotesize
\begin{listing}{1}
/*************************************************************************
* Процедура     : HandleCmdVis                                           *
*------------------------------------------------------------------------*
* Назначение    : Обработка визуализации полей                           *
* Входные       : hWnd - дескриптор окна, wmId - идентификатор команды   *
* Выходные      : Нет                                                    *
*------------------------------------------------------------------------*
* ОПИСАНИЕ ФУНКЦИОНИРОВАНИЯ:                                             *
* Инициирует рендеринг графиков через Wolfram (с проверкой готовности),  *
* управляет состоянием отображения и выводит ошибки при сбоях.           *
*************************************************************************/
void HandleCmdVis(HWND hWnd, int wmId) {
    if (wmId == ID_BTN_VIS_A || wmId == ID_BTN_VIS_B) {
        if (!g_IsCalculated) { ShowError(hWnd, WarnNoCalc); return; }
        bool isA = (wmId == ID_BTN_VIS_A); 
        HBITMAP& bmp = isA ? g_hPlotBitmapA : g_hPlotBitmapB;
        if (!bmp) { g_StatusMessage = isA ? L"⏳ Рендеринг Поля А..." : 
                                    L"⏳ Рендеринг Поля Б..."; g_ActiveView = 5; 
            InvalidateRect(hWnd, NULL, TRUE); UpdateWindow(hWnd); ErrCode err;
            if (!GenerateWolframPlot(isA, err)) {
                g_StatusMessage = L"❌ Ошибка рендеринга!"; g_ActiveView = 2;
                InvalidateRect(hWnd, NULL, TRUE); UpdateWindow(hWnd); ShowError(hWnd, err); 
                return;}
        }
        g_StatusMessage = isA ? L"✔ Визуализация Поля А готова!" : 
                                L"✔ Визуализация Поля Б готова!";
        g_ActiveView = isA ? 3 : 4;}
}
\end{listing}
\normalsize

\pagebreak

\section{Процедура HandleWmCommand}
\footnotesize
\begin{listing}{1}
/*************************************************************************
* Процедура     : HandleWmCommand                                        *
*------------------------------------------------------------------------*
* Назначение    : Диспетчеризация команд меню и кнопок                   *
* Входные       : hWnd - окно, wParam - параметры сообщения              *
* Выходные      : Нет                                                    *
*------------------------------------------------------------------------*
* ОПИСАНИЕ ФУНКЦИОНИРОВАНИЯ:                                             *
* Центральный обработчик сообщений WM_COMMAND, распределяющий логику      *
* в зависимости от нажатого элемента интерфейса.                         *
*************************************************************************/
void HandleWmCommand(HWND hWnd, WPARAM wParam) {
    int wmId = LOWORD(wParam);
    if (wmId == ID_BTN_IMPORT || wmId == ID_BTN_SHOW) 
        HandleCmdImportShow(hWnd, wmId);
    else if (wmId == ID_BTN_CALC || wmId == ID_BTN_EXPORT) 
        HandleCmdCalcExport(hWnd, wmId);
    else if (wmId == ID_BTN_VIS_A || wmId == ID_BTN_VIS_B) 
        HandleCmdVis(hWnd, wmId);
    InvalidateRect(hWnd, NULL, TRUE);
}
\end{listing}
\normalsize

\pagebreak

\section{Функция WndProc}
\footnotesize
\begin{listing}{1}
/*************************************************************************
* Функция       : WndProc                                                *
*------------------------------------------------------------------------*
* Назначение    : Основная процедура обработки событий окна              *
* Входные       : hWnd, msg, wParam, lParam (стандартные Win32)          *
* Выходные      : LRESULT - результат обработки события                  *
*------------------------------------------------------------------------*
* ОПИСАНИЕ ФУНКЦИОНИРОВАНИЯ:                                             *
* Главный цикл сообщений приложения, управляет жизненным циклом окна,    *
* отрисовкой и взаимодействием с пользователем.                          *
*************************************************************************/
LRESULT CALLBACK WndProc(HWND hWnd, UINT msg, WPARAM wParam, LPARAM lParam) {
    switch (msg) { case WM_CREATE: 
        CreateInterfaceButtons(hWnd, ((LPCREATESTRUCT)lParam)->hInstance); 
		break;
    case WM_COMMAND:  HandleWmCommand(hWnd, wParam); 
		break;
    case WM_PAINT: HandlePaint(hWnd); 
		break;
    case WM_DESTROY: if (g_hPlotBitmapA) DeleteObject(g_hPlotBitmapA);
        if (g_hPlotBitmapB) DeleteObject(g_hPlotBitmapB); PostQuitMessage(0); 
		break;
    default:  return DefWindowProc(hWnd, msg, wParam, lParam);}
    return 0;
}
\end{listing}
\normalsize

\pagebreak